\chapter{核心研究方法}

\section{基于双重修剪的分布式梯度下降算法}
为求解重构后的优化目标，本文提出了一种基于双重修剪机制的协同分布式梯度下降算法。算法具体执行流程如算法 \ref{alg:1} 所示（针对诚实智能体 $i\in\mathcal{V}\setminus\mathcal{A}$）。

\begin{algorithm}[H]
    \KwIn{$b$ 以及序列 ${\lambda[t]}_{t=1}^{\infty}$}
    将 $x_i[0]$ 设为约束集 $\mathcal{X}$ 中的任意值\;
    \For{$t = 1, \ldots$}{
        计算局部梯度 $f_i'(x_i[t-1])$ \;
        将 $(x_i[t-1], f_i'(x_i[t-1]))$ 发送给所有出边相连的邻居 \;
        $\mathcal{D}_i^x[t] \leftarrow$ 从入边相连的邻居收到的估计值集合 \;
        $\mathcal{D}_i^g[t] \leftarrow$ 收到的梯度集合 $\cup \{f_i'(x_i[t-1])\}$ \;
        $\widetilde{x}_i[t] \leftarrow \operatorname{Average}\bigl(Trim(\mathcal{D}_i^x[t]) \cup \{x_i[t-1]\}\bigr)$ \;
        $\widetilde{g}_i[t] \leftarrow \frac{1}{2}\bigl(\max\{Trim(\mathcal{D}_i^g[t])\} + \min\{Trim(\mathcal{D}_i^g[t])\}\bigr)$ \;
        计算 $x_i[t] \leftarrow P_{\mathcal{X}}[\widetilde{x}_i[t] - \lambda[t]\widetilde{g}_i[t]]$ \;
    }
    \caption{拜占庭容错梯度下降（诚实智能体 $i\in\mathcal{V}\setminus\mathcal{A}$ 的本地程序）}
    \label{alg:1}
    \noindent\textbf{函数 $Trim(\cdot)$：} \\
    \KwIn{多重集 $(\mathcal{D})$，其大小 $\ge 2b+1$} 
    将 $\mathcal{D}$ 中的元素按非降序排列后删除最小的 $b$ 个和最大的 $b$ 个元素 \\
    \KwOut{修剪后的集合 $\mathcal{D}$}
\end{algorithm}

在每次迭代中，最大拜占庭故障数 $b$ 和步长序列 $\{\lambda[t]\}_{t=1}^{\infty}$ 为已知参数。其中，步长满足单调不增 $\lambda[t] \ge \lambda[t+1]$、级数发散 $\sum \lambda[t] = \infty$ 和平方级数收敛 $\sum \lambda^2[t] < \infty$ 的标准凸优化条件。

诚实智能体 $i$ 在第 $t$ 次迭代时，将局部估计 $x_i[t-1]$ 和局部梯度 $f_i'(x_i[t-1])$ 传递给其出度邻居。随后，节点 $i$ 构建两个本地数据集合：$\mathcal{D}_i^x[t]$ 为从入度邻居收到的状态估计值集合，$\mathcal{D}_i^g[t]$ 为收到的梯度集合与自身局部梯度的并集。由于拜占庭节点可实施非一致性信息传输 \upcite{fischer1985impossibility}，不同诚实节点接收到的数据集合可能各不相同。若某入度邻居通信缺失，则以默认值处理，并自动将其视为拜占庭节点。

算法的核心在于引入了非线性截断函数 $Trim(\cdot)$。该函数接受规模不小于 $2b+1$ 的多重集合，对其进行非降序排列后，直接移除最小的和最大的各 $b$ 个元素。这一操作可在无需辨别节点身份的前提下，有效截断异常的极值。

在信息聚合阶段，状态与梯度的处理策略有所区别。

鲁棒状态估计 $\widetilde{x}_i[t]$ 取 $Trim(\mathcal{D}_i^x[t]) \cup \{x_i[t-1]\}$ 中所有元素的算术平均值；而鲁棒伪梯度 $\widetilde{g}_i[t]$ 取 $Trim(\mathcal{D}_i^g[t])$ 中最大值与最小值的平均数。随后，通过投影算子 $P_{\mathcal{X}}$ 将更新后的状态拉回至可行域 $\mathcal{X}$ 内避免状态偏离可行域，进而完成单次迭代：

\begin{equation}
    x_i[t] \leftarrow P_{\mathcal{X}}\bigl[\widetilde{x}_i[t] - \lambda[t]\widetilde{g}_i[t]\bigr].
    \label{eq:update}
\end{equation}

\section{算法收敛性分析}

在非独立同分布且去中心化的网络环境中，由于局部数据的异质性，系统在数学上难以区分诚实节点的极端观测值与拜占庭节点的对抗注入数据。因此，容错设计的核心转化为通过降维与鲁棒聚合函数，阻断异常数据值的扩散 \upcite{sundaram2019distributed}。

本文的分析表明，双重修剪机制能够人为构造出一个由诚实节点真实数据构成的几何边界。利用该凸包夹逼特性，可将含有对抗扰动的复杂动态系统等效映射为受限边界内的带误差梯度下降过程。在满足图连通性假设（详见附录 A.3）时，算法收敛性得到保证，详见定理\ref{thm:main}

\begin{theorem}
    设 $|\mathcal{N}_i\cap\mathcal{A}| \triangleq \phi_i$。定义参数边界为：
    \begin{align}
        \tilde{\gamma} &\triangleq \min_{i \in \mathcal{V} \setminus \mathcal{A}} (d_i + 1 - \phi_i - b), \\
        \tilde{\beta}  &\triangleq \min\left\{\frac{1}{2\max_{i \in \mathcal{V} \setminus \mathcal{A}}(d_i + 1 - \phi_i - b)},\ \frac{1}{n - \phi}\right\}.
    \end{align}
    则对所有诚实智能体 $i\in\mathcal{V}\setminus\mathcal{A}$，其状态渐近收敛于合法解集 $X(\widetilde{\beta}, \widetilde{\gamma})$，即 \\
    $\lim_{t\to\infty} \operatorname{Dist}\bigl(x_i[t], X(\widetilde{\beta}, \widetilde{\gamma})\bigr) = 0$。
    \label{thm:main}
\end{theorem}
